% ==STYLE-BEGIN==
\documentclass[UTF8,11pt]{ctexart}
\usepackage{amsmath,amssymb,amsthm}
\usepackage[a4paper,margin=2.4cm,top=2.3cm,bottom=2.3cm]{geometry}
\usepackage{fancyhdr}
\usepackage{enumitem}
\usepackage{hyperref}
\hypersetup{hidelinks}

% ---- 版式 ----
\linespread{1.15}
\setlength{\parskip}{0.35em}
\setlength{\parindent}{1.6em}
\setlist{nosep,leftmargin=2em}

% ---- 页眉 ----
\pagestyle{fancy}
\fancyhf{}
\fancyhead[L]{\small\coursename}
\fancyhead[R]{\small\thepage}
\renewcommand{\headrulewidth}{0.4pt}

% ---- 定理环境：中文用正体，标题加粗，按小节共用一条流水号 ----
\newtheoremstyle{cnthm}{0.7em}{0.7em}{\normalfont}{}{\bfseries}{}{0.5em}{}
\theoremstyle{cnthm}
\newtheorem{thm}{定理}[section]
\newtheorem{lem}[thm]{引理}
\newtheorem{prop}[thm]{命题}
\newtheorem{cor}[thm]{推论}
\newtheorem{defn}[thm]{定义}
\newtheorem{rem}[thm]{注}
\newtheorem{exa}[thm]{例}
\newtheorem{conj}[thm]{猜想}
\renewcommand{\proofname}{\bfseries 证明}
\providecommand{\coursename}{}
% ==STYLE-END==
\renewcommand{\coursename}{数学分析\,·\,1008}
\title{数学分析 1008}
\date{}
\begin{document}
\maketitle

\section{\ Recall}

\subsubsection*{1.\ 极限的定义}

对于
\[
\lim_{n\to+\infty}a_n=A,
\]
我们有
\[
\forall\varepsilon,\ \exists N,\ \forall n>N,\ \text{s.t.}\
|a_n-A|<\varepsilon .
\]
keypoint：放大不等式直至 $P(N)<\varepsilon$ 形式再反解。

\subsubsection*{2.\ 极限的性质}

\begin{itemize}
\item \textbf{有界性}\quad 有极限存在的数列必有界，即
\[
\text{if}\ \lim_{n\to+\infty}a_n=A\ \text{exists},\quad
\exists M,\ \forall n,\ |a_n|\leq M .
\]

\[
\lim_{n\to+\infty}a_n=0
\iff \lim_{n\to+\infty}|a_n|=0,
\]
直接使用极限语言易证。

\item \textbf{保号性}
\[
\begin{aligned}
&\text{if}\ \lim_{n\to+\infty}a_n=A\ \&\ A>0,
\quad\exists N,\ \forall n>N,\ a_n>\frac12A>0;\\
&\text{if}\ \lim_{n\to+\infty}a_n=A\ \&\ \exists N,\
\forall n>N,\ a_n>0,\quad A\geq0 .
\end{aligned}
\]
\end{itemize}

\subsubsection*{3.\ 极限的四则运算}

若
\[
\lim_{n\to+\infty}a_n=A,\qquad \lim_{n\to+\infty}b_n=B,
\]
则有
\[
\lim_{n\to+\infty}(a_n\pm\ \text{or}\ \times\ b_n)
=A\pm\ \text{or}\ \times\ B .
\]
如果 $B\neq0$，则有
\[
\lim_{n\to+\infty}\frac{a_n}{b_n}=\frac AB .
\]
保证 $\dfrac{a_n}{b_n}$ 是有定义的：根据极限保号性，舍去有限多项后
$b_n>0$，这保证舍去有限多项后 $\dfrac{a_n}{b_n}$ 都是有定义的。

\medskip

\textbf{e.g.}
\[
\lim_{n\to+\infty}
\frac{a_kn^k+a_{k-1}n^{k-1}+\cdots+a_1n+a_0}
{b_kn^k+b_{k-1}n^{k-1}+\cdots+b_1n+b_0},
\]
同除 $n^k$ 有
\[
\lim_{n\to+\infty}
\frac{a_k+\frac{a_{k-1}}{n}+\cdots+\frac{a_1}{n^{k-1}}
+\frac{a_0}{n^k}}
{b_k+\frac{b_{k-1}}{n}+\cdots+\frac{b_1}{n^{k-1}}
+\frac{b_0}{n^k}}
=\frac{a_k}{b_k},
\]
（上例中允许除数为零，规定 $\dfrac{a_k}{0}=+\infty$。）

\medskip

\textbf{e.g.}\quad 设 $\lim\limits_{n\to+\infty}a_n=A$，
$\lim\limits_{n\to+\infty}b_n=B$，$A>B$，则 $\exists N$，
$\forall n>N$ 有 $a_n>b_n$（保序性）。

考虑
\[
\lim_{n\to+\infty}(a_n-b_n)=A-B>0,
\]
据保号性 $\exists N$，$\forall n>N$，有 $a_n-b_n>0$，
即 $a_n>b_n$。

反过来也成立：设 $\lim\limits_{n\to+\infty}a_n=A$，
$\lim\limits_{n\to+\infty}b_n=B$，且 $\exists N$，$\forall n>N$，
$a_n\geq b_n$，则有 $A\geq B$。

\subsubsection*{4.\ 夹迫准则}

若 $\exists N$ s.t. $\forall n>N$，有
$a_n\leq c_n\leq b_n$ 且
\[
\lim_{n\to\infty}a_n=\lim_{n\to\infty}b_n=A,
\]
则
\[
\lim_{n\to\infty}c_n=A .
\]

\medskip

\textbf{e.g.}（已知，见上题 (2)）
\[
\lim_{n\to+\infty}\sqrt[n]{a}=1\quad(a>0).
\]
下面试证
\[
\lim_{n\to+\infty}\sqrt[n]{n}=1 .
\]

\textbf{证明：}\quad
\[
\sqrt[n]{n}>1\ \Longrightarrow\
\sqrt[n]{n}-1\triangleq h_n
\iff \sqrt[n]{n}=1+h_n,
\]
\[
\Longrightarrow\ n=(1+h_n)^n
=1+nh_n+\frac{n(n-1)}{2}h_n^2+\cdots
\]
其中，先尝试保留一次项：$n\geq1+nh_n\Rightarrow h_n\leq\dfrac{n-1}{n}$，次数不够（只得到 $O(1)$ 的界，推不出 $h_n\to0$）；故保留二次项（其余各项非负，舍去即得下界），即
\[
n\geq\frac{n(n-1)}{2}h_n^2 .
\]
易得
\[
\sqrt{\frac{2}{n-1}}\geq h_n>0,
\]
由夹迫准则即可得到结果。

\section{\ 单调有界数列的极限}

\begin{defn}
设 $\{a_n\}$ 满足 $\forall n,m$，$n<m$ 有
$a_n\leq a_m$，则称 $\{a_n\}$ 单调上升；若 $\forall n,m$，
$n<m$ 有 $a_n<a_m$，则称严格单调。
\end{defn}

\begin{thm}
设 $\{a_n\}$ 单调递增且 $\{a_n\}$ 有上界，则
\[
\lim_{n\to+\infty}a_n
\]
存在。
\end{thm}

\subsubsection*{分析}

\[
S=\{b\mid b\ \text{是}\ \{a_n\}\ \text{的上界}\},\qquad
T\triangleq\mathbb R-S .
\]
我们需要证明：

\begin{enumerate}
\item $T\mid S$ 是分划；
\item $\mathbb R$ 是 Dedikind 连续的。

以上两条我们知道存在界点 $A$。
\end{enumerate}

\noindent 3.\ 界点 $A\in S$。

\subsubsection*{证明}

易知
\[
S\cup T=\mathbb R,\qquad S\cap T=\varnothing .
\]
要证明它是分划，只需证明
\[
\forall a\in T,\ \forall b\in S,\quad a<b .
\]
\[
\begin{aligned}
&a\in T\iff a\ \text{不是上界}\iff\exists m,\ \text{s.t.}\ a_m>a;\\
&b\in S\iff b\ \text{是上界}\iff\forall n,\ \text{s.t.}\ a_n\leq b .
\end{aligned}
\]

特别地，取 $n=m$，有 $a<a_m\leq b$。可知 $T\mid S$ 确实是分划。
而分划 $T\mid S$ 一定有界点 $A$。

下证界点 $A\in S$。反设 $A$ 不为上界，则 $\exists m$，
$a_m>A$。可知 $\exists B$ s.t. $a_m>B>A$。

由 $B>A$ 及 $A$ 是界点知 $B\in S$，即 $B$ 是上界；
但 $B<a_m$ 与 $B$ 是上界矛盾。假设不成立，即 $A\in S$——
$A$ 是上界中最小的，即上确界。

下证明
\[
\lim_{n\to+\infty}a_n=A .
\]
由极限定义可知，上式等价于
\[
|a_n-A|<\varepsilon
\iff A-\varepsilon<a_n<A+\varepsilon .
\]
其中 RHS 自然成立，因为 $A$ 已经是上界了。
$A-\varepsilon\in T$，所以它不是上界，即
\[
\exists N,\ \text{s.t.}\ a_n\geq a_N\geq A-\varepsilon .
\]
可得
\[
\lim_{n\to+\infty}a_n=A .
\]

\end{document}
